\section{First-Order Phase Transitions and Phase Diagram}
For a given system, a \textcolor{blue}{phase} is a region of the space of natural
variables in which the thermodynamic potential is analytic: infinitely
differentiable, with no kinks or divergences. A \textcolor{blue}{phase transition}
is a point where it fails to be analytic. Strictly, such singularities appear only
in the thermodynamic limit $N\to\infty$, since the partition function of a finite
system is analytic.

We usually study phases at fixed particle numbers, pressure and temperature, for
two reasons. Experimentally, these are easy to control. Structurally, the phases
inside such a system exchange energy, volume and particles with each other, so at
equilibrium $(\{\mu_i\},P,T)$ must be equal in all phases. The natural choice of
potential is then the Gibbs free energy $G(T,P,\{N_i\})$. Because $G$ is extensive,
\begin{align*}
    G(T,P,\{N_i\}) = N\,g(T,P,\{x_i\}), \qquad x_i = N_i/N,
\end{align*}
so the total size $N$ does not affect the analyticity of $G$; only $g$ does.
What matters is the composition $\{x_i\}$, of which $C-1$ are independent for $C$ component. A phase is therefore a region of the intensive variables $(T,P,\{x_i\})$,
which span a $(C-1)+2=C+1$ dimensional space. Collecting every phase in this
space gives the \textcolor{blue}{phase diagram}.

At equilibrium, $g$ is continuous, so a phase transition is never a jump in the
potential itself. Adding $g$ as one more axis to the $(C+1)$-dimensional space
$(T,P,\{x_i\})$ gives a hypersurface $g(T,P,\{x_i\})$. A
\textcolor{blue}{first-order transition} is a crease (kink) of this hypersurface:
$g$ is continuous across it, but its first derivatives jump,
\begin{align*}
    s = -\frac{\partial g}{\partial T}, \qquad v = \frac{\partial g}{\partial P},
\end{align*}
the entropy and volume per particle. Equivalently, $S=-\partial G/\partial T$ and
$V=\partial G/\partial P$ are discontinuous, which gives a latent heat
$L = T\Delta S$ and a volume change $\Delta V$. Melting and boiling are
examples. On the transition locus the coexisting phases share the same
$(\{\mu_i\},P,T)$ but have different $s$, $v$ and composition, which is why the
transition appears as a boundary of a coexistence region in the phase diagram.

Transitions of higher order also exist: the order is the order of the lowest
derivative of $g$ that is discontinuous or singular, and for these the
hypersurface has no kink and the singularity appears in the curvature. They
include the continuous transitions, such as the ferromagnetic Curie point, where
response functions like the heat capacity diverge. We defer them to a later
section.

\subsection{Gibbs Phase Rule}

    Suppose $\mathcal{P}$ phases of a system with $C$ components coexist. At
    equilibrium all phases share the same intensive variables
    $(T,P,\mu_1,\dots,\mu_C)$, which gives $C+2$ variables. Each phase has its
    own Gibbs--Duhem relation,
    \begin{align*}
        S_\alpha\,dT - V_\alpha\,dP + \sum_i N_{i,\alpha}\,d\mu_i = 0,
        \qquad \alpha = 1,\dots,\mathcal{P},
    \end{align*}
    which gives $\mathcal{P}$ independent constraints on them. The number of
    degrees of freedom $F$ is the number of variables minus the number of
    constraints,
    \begin{align}
        \boxed{F+\mathcal{P}=C+2.}
    \end{align}
    This is the \textcolor{blue}{Gibbs phase rule}. It counts only intensive
    degrees of freedom: the amount of each phase remains free and is not counted. However, this is only applicable if $(T,P,\mu_1,\dots,\mu_C)$ are the only intensive variables. We will encounter system that depends on the magnetic field $\vec B$ of itself, which add another intensive variable to the count.

\subsection{Clapeyron Equation and Fixed-Composition ($P$--$T$) Phase Diagram}
No matter how many components the system has, temperature and pressure always
participate. Consider the $P$--$T$ plane at fixed composition $\{x_i\}$ within the
$(C+1)$-dimensional phase diagram. In this slice only $(P,T)$ can vary, so there
are $2$ degrees of freedom. Two phases $1$ and $2$ coexist when every chemical
potential matches,
\begin{equation*}
    \mu_i^{(1)}(T,P) = \mu_i^{(2)}(T,P), \qquad i=1,\dots,C,
\end{equation*}
which relates $P$ and $T$ and reduces the degrees of freedom from $2$ to $1$.
The coexistence locus is therefore a curve in the plane. Moving along it, the two
phases keep equal chemical potentials, so $d\mu_i^{(1)}=d\mu_i^{(2)}$ and, by
the Gibbs--Duhem relation at fixed composition,
\begin{equation*}
    -s_1\,dT + v_1\,dP = -s_2\,dT + v_2\,dP ,
\end{equation*}
which gives the \textcolor{blue}{Clapeyron equation}
\begin{equation}
    \boxed{\frac{dP}{dT}=\frac{\Delta s}{\Delta v}}, \qquad
    \Delta s = s_2-s_1,\;\; \Delta v = v_2-v_1 .
\end{equation}
At coexistence the transition proceeds reversibly at fixed $T$ and $P$, so the
latent heat per particle is $\ell = T\Delta s$, and
\begin{equation}
    \boxed{\frac{dP}{dT}=\frac{\ell}{T\,\Delta v}.}
\end{equation}

Now suppose a third phase $3$ also coexists. Then
\begin{equation*}
    \mu_i^{(1)}(T,P)=\mu_i^{(2)}(T,P)=\mu_i^{(3)}(T,P), \qquad i=1,\dots,C .
\end{equation*}
Each pair of phases has its own coexistence curve, $\mu^{(1)}=\mu^{(2)}$,
$\mu^{(2)}=\mu^{(3)}$ and $\mu^{(1)}=\mu^{(3)}$. A point lying on the first two
satisfies $\mu^{(1)}=\mu^{(2)}=\mu^{(3)}$, so it lies on the third as well:
only two of the three conditions are independent. Two independent equations in the
two unknowns $(T,P)$ leave no degrees of freedom, so three phases can coexist
only at an isolated point, the \textcolor{blue}{triple point}, where the three
two-phase curves, the \textcolor{blue}{triple curves} intersect.

\subsection{Unary System}
A \textcolor{blue}{unary system} has only one component, $C=1$. This is a special case of the
fixed-composition slice, since the composition is always $x=1$. We discuss it
because it is the building block of higher-dimensional phase diagrams and
introduces some terminology. For $C=1$ the phase rule reads
$F = 3-\mathcal{P}$, so the number of coexisting phases $\mathcal{P}$ fixes the
number of free variables.

Within a single phase ($\mathcal{P}=1$, $F=2$), $P$ and $T$ can both be varied
freely. The system is said to be in \textcolor{blue}{bivariant equilibrium}.

When two phases coexist ($\mathcal{P}=2$, $F=1$), only one of $P$ and $T$ can be
varied freely; the other is then fixed by the Clapeyron relation $P=P(T)$. The
system is said to be in \textcolor{blue}{univariant equilibrium}, and the locus
of two-phase coexistence is called a \textcolor{blue}{univariant curve}.

When three phases coexist ($\mathcal{P}=3$, $F=0$), no variable can be changed
freely. The system is said to be in \textcolor{blue}{invariant equilibrium}, and
the point where this happens is called an \textcolor{blue}{invariant point}
(the triple point).

    